\chapter{CRESCERE NELL'ACQUA BOLLENTE}\label{crescere-nell-acqua-bollente}

\hfill\break

Quelli più giovani di me mi pongono spesso domande su quanto \emph{non}
accadde. Perché non siete andati nel panico? Perché \emph{nessuno} è
andato nel panico? Perché non c'è stato nessun saccheggio, nessuna
rivolta? Perché la tua generazione ha accettato supinamente, perché
siete scivolati nello Spin senza neanche un cenno di protesta?

A volte rispondo: «Ma accaddero cose davvero terribili.»

A volte rispondo: «Ma noi non capivamo. E cosa avremmo potuto farci?»

E a volte cito la parabola della rana: «Getta una rana nell'acqua
bollente, vedrai che salta fuori. Getta una rana in un contenitore con
dell'acqua tiepida e piacevole, metti la fiamma bassa e la rana morirà
prima di rendersi conto che è nei guai.»

L'annientamento delle stelle non fu lento o impercettibile, ma nemmeno
disastroso nell'immediato, almeno per la maggior parte di noi. Un
astronomo o uno stratega della difesa, chi lavorava nelle
telecomunicazioni o nell'ambito aerospaziale aveva forse passato i primi
giorni dello Spin in uno stato di terrore assoluto. Ma per uno che
guidava gli autobus o rivoltava hamburger doveva essere stato più o meno
come essere nell'acqua tiepida.

I media lo chiamarono ``l'Evento di ottobre'' (divenne ``lo Spin'' solo
qualche anno dopo) e il primo effetto evidente fu la distruzione su
larga scala dell'industria multimiliardaria dei satelliti orbitali. Con
la perdita dei satelliti, restammo senza gran parte delle reti
televisive, a partire da tutte quelle trasmesse proprio via satellite;
le linee telefoniche a lunga distanza divennero inaffidabili e i
localizzatori GPS inutili; il World Wide Web fu annientato e molta della
più sofisticata e moderna tecnologia militare divenne obsoleta. La
sorveglianza e la ricognizione mondiali furono ridotte e i meteorologi
locali ripresero a disegnare le isobare sulle cartine degli Stati Uniti
anziché scivolare sulle immagini fornite dai satelliti ed elaborate al
computer. I ripetuti tentativi di contattare la Stazione Spaziale
Internazionale furono vani. I lanci commerciali programmati a Canaveral
(e a Baikonur e Kourou) furono rinviati a tempo indefinito.

Nel lungo periodo tutto questo significava un mare di guai per GE
Americom, AT\&T, COMSAT e Hughes Communications, fra molte altre.

E \emph{accaddero} cose davvero terribili come conseguenza di quella
sera, anche se la maggior parte di esse furono nascoste dai black-out
dei media. Le notizie viaggiarono come sussurri lungo i cavi della fibra
ottica anziché rimbalzare nello spazio orbitale: passò quasi una
settimana prima di sapere che un missile pachistano Hatf V con una
testata nucleare, lanciato per errore o per un calcolo sbagliato nei
primi momenti di confusione dell'Evento, era finito fuori rotta e aveva
vaporizzato una valle coltivata dell'Hindu Kush. Era il primo
dispositivo nucleare esploso in guerra dal 1945 e, vista la tragicità
dell'evento e data la paranoia globale innescata dalla perdita delle
telecomunicazioni, fummo fortunati che accadde una volta sola. Secondo
alcuni resoconti stavamo quasi per perdere Teheran, Tel Aviv e
Pyongyang.

\separatorefiori

Rassicurato dal sorgere del Sole, dormii dall'alba a mezzogiorno. Quando
mi alzai, trovai mia madre in soggiorno. Indossava ancora la vestaglia
trapuntata e fissava pensierosa lo schermo della televisione. Quando le
chiesi se avesse fatto colazione mi rispose di no. E allora preparai il
pranzo per entrambi.

Quell'autunno avrebbe compiuto quarantacinque anni. Se mi avessero
chiesto di scegliere una parola per descriverla sarebbe stata
``solida''. Si arrabbiava di rado e l'unica volta in vita mia che
l'avevo vista piangere era stata la sera in cui venne a casa la polizia
(ai tempi di Sacramento) per dirle che mio padre era morto
sull'Autostrada 80 vicino a Vacaville, mentre tornava a casa da un
viaggio di lavoro. Credo che facesse attenzione a mostrarmi solo questo
aspetto di sé. Ma ce n'erano altri. C'era un ritratto su una mensola
dell'étagère nel soggiorno, scattato anni prima che nascessi, di una
donna così elegante, bella e sicura di sé davanti alla macchina
fotografica che ero rimasto sbigottito quando mi aveva detto che era una
sua foto.

Era evidente che non le piaceva quello che stava ascoltando alla TV.
Un'emittente locale trasmetteva notiziari non stop, riprendendo i
resoconti forniti dai radioamatori, dalle trasmissioni a onda corta e
diffondendo i comunicati poco chiari diramati dal governo federale che
invitavano a mantenere la calma. «Tyler,» mi disse facendo segno di
sedermi, «è difficile da spiegare. È successo qualcosa ieri sera\ldots»

«Lo so» risposi. «Ne ho sentito parlare prima di andare a letto.»

«Lo sapevi? E non mi hai svegliato?»

«Non sapevo se\ldots»

Ma il suo fastidio scemò presto così come era arrivato. «No,» mi
interruppe, «non importa, Ty. Credo di non essermi persa niente. È
buffo\ldots{} mi sento come se stessi \emph{ancora} dormendo.»

«Sono solo stelle» replicai stupidamente.

«Stelle e Luna» mi corresse lei. «Non hai sentito della Luna? In tutto
il mondo non si vedono né le stelle né la Luna.»

\separatorefiori

Certamente la Luna era un indizio.

Rimasi ancora un po' seduto con mia madre, poi la lasciai ancora
immobile davanti alla TV («Stavolta torna prima che faccia buio» mi
ammonì seria) e me ne andai a piedi alla Grande Casa. Bussai alla porta
sul retro, quella che usavano il cuoco e la domestica diurna, anche se i
Lawton badavano bene a non chiamarla mai ``entrata di servizio''. Era
anche la porta da cui entrava mia madre per gestire le questioni
domestiche dei Lawton durante la settimana.

La signora Lawton, la madre dei gemelli, mi fece entrare, mi guardò con
aria assente e mi fece cenno di andare di sopra. Diane stava ancora
dormendo, la porta della sua camera era chiusa. Jason non aveva dormito
e sembrava non ne avesse alcuna intenzione. Lo trovai in camera sua che
monitorava una radio a onde corte.

La stanza di Jason era come una grotta delle meraviglie di Aladino, ne
desideravo il contenuto ma avevo smesso di sperare di possederlo: un
computer dotato di connessione internet ultraveloce e una televisione di
seconda mano due volte più grande di quella che ornava il soggiorno di
casa mia. Pensando che potesse non avere sentito la notizia, lo
informai: «La Luna è sparita.»

«Interessante, no?» Jase si alzò in piedi e si stiracchiò, passandosi le
dita fra i capelli spettinati. Non si cambiava dalla sera prima. Era una
distrazione insolita. Jason, anche se lo era per davvero, non si era mai
atteggiato a genio in mia presenza - cioè, non si comportava come i geni
che avevo visto nei film; non era strabico, non balbettava né scriveva
equazioni algebriche sui muri. Eppure quel giorno sembrava distratto.
«La Luna \emph{non} è sparita, ovviamente - come potrebbe? Secondo la
radio vengono rilevate le solite maree sulla costa atlantica. Perciò la
Luna è ancora al suo posto. E se la Luna è ancora lì, lo sono anche le
stelle.»

«E allora perché non le vediamo?»

Mi lanciò uno sguardo seccato. «E che ne so io? Dico solo che almeno in
parte è un fenomeno \emph{ottico}.»

«Guarda fuori dalla finestra, Jase. Il Sole splende. Che genere di
illusione ottica lascia passare la luce del Sole ma nasconde la Luna e
le stelle?»

«Te lo ripeto, che ne so? Ma qual è l'alternativa, Tyler? Qualcuno ha
messo la Luna e le stelle in un sacco ed è scappato portandosele via?»

``No'' pensai. Era la Terra a essere nel sacco, per qualche motivo che
neanche Jason sapeva interpretare.

«Comunque, bella ipotesi quella sul Sole. Non una barriera ottica ma
\emph{un filtro} ottico. Interessante\ldots»

«E quindi chi ce l'ha messo?»

«E che ne\ldots?» Scosse la testa, irritato. «Salti a conclusioni troppo
affrettate. Chi lo dice che ce l'ha messo \emph{qualcuno}? Potrebbe
essere un evento naturale di quelli che capitano una volta ogni milione
di anni, come quello dell'inversione dei poli magnetici. È un grosso
salto supporre che dietro ci sia un'intelligenza che controlla tutto.»

«Ma potrebbe essere vero.»

«Tante cose potrebbero essere vere.»

Ero stato già abbastanza preso in giro, pur se in maniera lieve, per le
mie letture di fantascienza da essere restio a pronunciare la parola
``alieni''. Ma ovviamente fu la prima cosa che mi venne in mente. A me e
a molta altra gente. E anche Jason dovette ammettere che l'idea di un
intervento degli extraterrestri era diventata molto più plausibile nel
corso delle ultime ventiquattro ore.

«Ammettendo che sia vero,» dissi, «bisogna chiedersi perché dovrebbero
farlo.»

«Ci sono solo due ragioni plausibili. Nasconderci \emph{qualcosa}. O
nasconderci da \emph{qualcosa}.»

«Cosa ne pensa tuo padre?»

«Non gliel'ho chiesto. È tutto il giorno che è al telefono. Forse sta
cercando di piazzare il prima possibile le sue azioni GTE.» Era una
battuta, e non ero sicuro di cosa volesse intendere, ma fu anche il mio
primo indizio su ciò che la perdita dell'accesso all'orbita potesse
significare per l'industria aerospaziale in generale e per la famiglia
Lawton in particolare. «Stanotte non ho dormito» ammise Jase. «Temevo di
perdermi qualcosa. A volte invidio mia sorella. Sai com'è lei,
``svegliami quando ci capiscono qualcosa.''»

Mi irritai per quella velata offesa nei confronti di Diane. «Nemmeno lei
ha dormito.»

«Ah, davvero? E \emph{tu} che ne sai?»

Beccato. «Abbiamo parlato per un po' al telefono\ldots»

«Ti ha chiamato lei?»

«Sì, verso l'alba.»

«Cavolo, Tyler, stai arrossendo.»

«No.»

«Sì.»

Mi salvò un brusco colpo alla porta: era E.D. Lawton, anche lui sembrava
non aver dormito granché.

Il padre di Jason era una presenza intimidatoria. Grosso, spalle larghe,
difficile da accontentare e facilmente irascibile, durante il fine
settimana girava per casa come un uragano, tutto tuoni e fulmini. Mia
madre una volta aveva detto: «E.D. non è il tipo di persona da cui
vorresti ricevere attenzioni. Non ho mai capito perché Carol lo abbia
sposato.»

Non era proprio il classico uomo d'affari che si è fatto da solo - suo
nonno, fondatore in pensione di uno studio legale a San Francisco ben
avviato, aveva finanziato gran parte dei primi investimenti di E.D. - ma
lui si era costruito un'azienda redditizia in strumentazione d'alta
quota e tecnologia aerostatica, e aveva fatto tutto da sé con gran
fatica, senza veri agganci nell'industria, almeno quando aveva
cominciato.

Entrò con un'espressione torva nella camera di Jason. Posò lo sguardo su
di me e lo distolse in un lampo. «Mi spiace, Tyler, ma ora devi andare a
casa. Devo parlare un po' con Jason.»

Jase non protestò e neanch'io morivo dalla voglia di rimanere. E così mi
strinsi nella mia giacca di stoffa e uscii dalla porta sul retro.
Trascorsi il resto del pomeriggio vicino al ruscello, a far rimbalzare
pietre e a guardare gli scoiattoli in cerca di cibo in vista
dell'inverno imminente.

\separatorefiori

Il Sole, la Luna e le stelle.

Negli anni a seguire i bambini crebbero senza aver mai visto la Luna con
i loro occhi; persone di soli cinque o sei anni più giovani di me
divennero adulte conoscendo le stelle soprattutto da vecchi film e
attraverso una manciata di cliché sempre più inappropriati. Ricordo che
avevo circa trent'anni, quando suonai \emph{Corcovado} per una ragazza
più giovane. Era un pezzo che Antonio Carlos Jobim aveva composto nel
ventesimo secolo e che iniziava con le parole ``\emph{Notti silenti di
	stelle silenti}''. Sentendole, lei mi chiese, con sguardo genuinamente
sorpreso: «Erano \emph{rumorose} le stelle?»

Avevamo perso qualcosa di più impercettibile di qualche luce nel cielo.
Avevamo perso il senso di affidabilità di un luogo. La Terra è rotonda,
la Luna gira intorno alla Terra, la Terra gira intorno al Sole: erano
quelle le nozioni di cosmologia che la gente possedeva o si faceva
bastare, e dubito che dopo le scuole superiori più di una persona su
cento avesse ripensato a quelle cose. Ma rimanemmo tutti ugualmente
sconcertati, quando ne fummo privati.

Non ricevemmo alcun comunicato ufficiale sulla situazione del Sole fino
alla seconda settimana dall'Evento di ottobre.

Il Sole pareva muoversi nel solito modo eterno e prevedibile. Sorgeva e
tramontava secondo le normali effemeridi, i giorni si accorciavano
secondo la loro naturale precessione; non c'era niente che facesse
pensare a un'emergenza solare. Sulla Terra molte cose, inclusa la vita
stessa, erano sempre dipese dalla natura e dalla quantità di radiazione
solare che investiva la superficie del pianeta, e per molti aspetti
questo non era cambiato. Il Sole che vedevamo a occhio nudo lasciava
intendere che quella era la stella gialla di classe G di sempre, quella
verso cui avevamo strizzato gli occhi tutta la vita.

Mancavano invece le macchie solari, le protuberanze e i brillamenti.

Il Sole è un oggetto tumultuoso, violento. Bolle, ribolle, vibra come
una campana con immense energie; inonda il sistema solare con un flusso
di particelle cariche che ci ucciderebbero se non fossimo protetti dal
campo magnetico della Terra. Ma dall'Evento di ottobre, annunciarono gli
astronomi, il Sole era diventato una sfera geometricamente perfetta di
una lucentezza ferma, uniforme e pura. Giunse notizia dal Nord che
l'aurora boreale, frutto dell'interazione del nostro campo magnetico con
tutte quelle particelle solari cariche, aveva chiuso i battenti come una
brutta commedia di Broadway.

Altre perdite nel nuovo cielo notturno: niente stelle cadenti. La Terra
ogni anno aggregava trentaseimila tonnellate di polvere spaziale, di cui
la stragrande maggioranza incenerita dall'attrito atmosferico. Ma ora
non più: nessun meteorite rilevabile era penetrato nell'atmosfera
durante le prime settimane dall'Evento di ottobre, neppure quelle
microscopiche particelle chiamate Brownlee. C'era, in campo astrofisico,
un silenzio assordante.

Neanche Jason riusciva a darne una spiegazione.

\separatorefiori

E così quel \emph{sole} non era il Sole; ma continuava a brillare,
contraffatto o no, e mentre passavano i giorni, giorni su giorni
stratificati e accatastati, il disorientamento aumentava ma diminuiva la
sensazione di urgenza pubblica. (L'acqua non bolliva, era solo calda.)

Ma che ricca fonte di \emph{chiacchiericcio} erano diventati il mistero
celeste e le sue conseguenze immediate: il crollo delle
telecomunicazioni; le guerre estere non più monitorate e raccontate via
satellite; le bombe intelligenti guidate dal GPS rese irrimediabilmente
stupide; la corsa all'oro della fibra ottica. Da Washington venivano
rilasciate dichiarazioni con deprimente regolarità: ``Non abbiamo ancora
alcuna prova di intento ostile da parte di nessuna nazione o agenzia'' e
``Le menti migliori della nostra generazione stanno lavorando per
capire, spiegare e, in ultima analisi, indicare i potenziali effetti
negativi di questa coltre che ha oscurato la nostra visione
dell'universo.'' Frasi rassicuranti ma senza senso da parte di
un'amministrazione che sperava ancora di identificare un nemico,
terrestre o no, capace di una simile azione. Un nemico che si stava
rivelando, tuttavia, ostinato nel suo essere sfuggevole. La gente
cominciò a parlare di ``un'ipotetica intelligenza controllante''.
Incapaci di vedere oltre le mura della nostra prigione, eravamo ridotti
a mapparne i bordi e gli angoli.

Dopo l'Evento, Jason rimase rintanato nella sua stanza per quasi un
mese. Durante quel periodo non gli parlai mai di persona, lo vedevo di
sfuggita solo quando il minibus della Rice Academy passava a prendere i
gemelli. Tuttavia, Diane mi chiamava al cellulare quasi ogni sera, di
solito verso le dieci o le undici, quando entrambi potevamo avere un po'
di privacy. E io facevo tesoro delle sue chiamate, per motivi che non
ero ancora pronto ad ammettere a me stesso.

«Jason è proprio incazzato» mi disse una sera. «Dice che se non sappiamo
se il \emph{sole} è il Sole, allora non sappiamo proprio niente.»

«Forse ha ragione.»

«Ma è quasi una roba religiosa per Jase. Ha sempre adorato le carte - lo
sapevi, Tyler? Anche quando era piccolissimo, capiva come funzionava una
carta. Gli piaceva sapere dove si trovava. ``Dà senso alle cose'' diceva
sempre. Cavolo, mi piaceva da morire sentirlo parlare di carte. Credo
sia per questo che ora dà di matto più di chiunque altro. Niente è dove
dovrebbe essere. Ha perso la sua carta.»

Naturalmente c'erano già gli indizi giusti. Prima della fine della
settimana l'esercito aveva cominciato a raccogliere i detriti dei
satelliti caduti - satelliti che erano stati stabilmente in orbita fino
a quella sera di ottobre, per essere poi ricacciati in blocco sulla
Terra prima dell'alba, alcuni lasciando rottami che avrebbero assunto
valore di prova. Ci volle però del tempo perché quell'informazione
raggiungesse la famiglia di E.D. Lawton, pure ben ammanicata.

\separatorefiori

Il nostro primo inverno di notti buie fu claustrofobico e strano. La
neve arrivò in anticipo: vivevamo a poche ore da Washington D.C. ma a
Natale sembrava di stare nel Vermont. Le notizie restavano inquietanti.
Un trattato di pace tra India e Pakistan, fragile e negoziato
frettolosamente, vacillava ora sì ora no verso la guerra; il progetto di
decontaminazione sostenuto dalle Nazioni Unite nell'Hindu Kush era già
costato decine di vite oltre alle perdite iniziali. Nell'Africa
settentrionale covavano i conflitti locali mentre gli eserciti del mondo
industriale si ritiravano per riorganizzarsi. I prezzi del petrolio
salivano alle stelle. In casa tenevamo il termostato un paio di gradi
sotto la temperatura confortevole, almeno finché i giorni non
cominciarono ad allungarsi (``\emph{quando il Sole sorse, al primo verso
	della quaglia}'').

Bisogna tuttavia riconoscere che, di fronte a quelle minacce sconosciute
e inafferrabili, la specie umana riuscì a non scatenare una guerra
mondiale in piena regola. Ci adattammo e continuammo con la nostra vita,
e in primavera la gente cominciò a parlare della ``nuova normalità''.
Nel lungo periodo, ormai si era capito, avremmo dovuto pagare un prezzo
più alto per ciò che era successo al pianeta, qualsiasi cosa
fosse\ldots{} ma nel lungo periodo saremo tutti morti, come si suol
dire.

Notai il cambiamento in mia madre. Il passare del tempo la calmò e il
caldo, quando finalmente giunse, le strappò via un po' della tensione
dal viso. E notai il cambiamento anche in Jason, che uscì dal suo
rifugio meditativo. Mi preoccupavo, però, di Diane, che si rifiutava
categoricamente di parlare delle stelle e di recente aveva iniziato a
chiedermi se credessi in Dio e se pensassi che Dio fosse responsabile di
ciò che era avvenuto in ottobre.

«Non saprei» le rispondevo. La mia famiglia non era praticante.
Sinceramente l'argomento mi rendeva un po' nervoso.

\separatorefiori

Quell'estate noi tre andammo in bicicletta al centro commerciale Fairway
Mall per l'ultima volta.

Avevamo fatto quel tragitto cento, mille altre volte. I gemelli ormai
erano un po' grandicelli per queste cose ma nei sette anni vissuti nella
tenuta della Grande Casa era diventato un rituale, l'appuntamento
inderogabile dei nostri sabati estivi. Saltavamo l'incontro nei fine
settimana piovosi o di caldo soffocante, ma quando il tempo era bello
eravamo attratti come da una mano invisibile verso il nostro punto di
ritrovo alla fine del lungo viale dei Lawton.

Quel giorno l'aria era dolce e frizzante e la luce infondeva in tutto
ciò che toccava un intenso caldo organico. Era come se il clima volesse
rassicurarci: il mondo naturale se la cavava bene, grazie, quasi dieci
mesi dopo l'Evento\ldots{} anche se ora eravamo (come di tanto in tanto
diceva Jase) un pianeta \emph{coltivato}, un giardino curato da forze
sconosciute anziché un appezzamento di foresta cosmica.

Jason montava una mountain bike costosa, Diane l'equivalente femminile
ma meno appariscente. La mia bici era un catorcio di seconda mano che
mia madre aveva comprato in un mercatino dell'usato. Non era importante.
Ciò che invece contava era il forte odore di pino dell'aria e il
susseguirsi delle ore libere che ci attendevano. Io lo percepivo, Diane
lo percepiva e credo che anche Jason lo percepisse, nonostante sembrasse
distratto e persino un po' imbarazzato quando quella mattina montammo in
sella. Imputai il suo atteggiamento allo stress o (era agosto) alla
prospettiva di un altro anno scolastico. Alla Rice, Jase era in una
sezione scolastica accelerata, una scuola ad alta pressione. L'anno
prima aveva superato facilmente i corsi di matematica e fisica - che
avrebbe potuto tenere lui - ma nel semestre successivo era stato
iscritto, per ottenere crediti, a latino. «Non è nemmeno una lingua
viva» aveva detto. «Chi diavolo legge il latino, a parte i classicisti?
È come imparare il Fortran. Tutti i testi importanti sono stati tradotti
molto tempo fa. Mi rende una persona migliore leggere Cicerone in lingua
originale? Cicerone, andiamo! L'Alan Dershowitz della Repubblica
Romana!»

Non presi niente di ciò che disse troppo sul serio. Una delle cose che
ci piaceva fare durante queste corse era praticare l'arte di lamentarsi.
(Non avevo idea di chi fosse Alan Dershowitz; forse un ragazzino della
scuola di Jason.) Ma quel giorno il suo umore era volubile,
imprevedibile. Si sollevò sui pedali e corse per un breve tratto davanti
a noi.

La strada per il centro commerciale si snodava oltre appezzamenti
fittamente alberati e case pastello con giardini curati e irrigatori
nascosti che segnavano l'aria mattutina con arcobaleni. La luce solare
poteva essere finta, filtrata, ma si divideva ancora nei diversi colori
quando passava attraverso l'acqua in caduta e sembrava ancora una
benedizione quando rotolavamo da sotto le querce ombreggianti sul
marciapiede bianco splendente.

Dopo dieci o quindici minuti di pedalata leggera si profilava davanti a
noi la sommità della Bantam Hill Road - ultimo ostacolo e punto di
riferimento importante sulla strada per il centro commerciale. La Bantam
Hill Road era ripida ma dall'altra parte era una lunga discesa dolce
verso l'area di parcheggio del centro commerciale. Jase era già a un
quarto della salita. Diane mi lanciò un'occhiata maliziosa.

«A chi arriva primo!» esclamò.

Era sconcertante. Il compleanno dei gemelli era a giugno. Il mio a
ottobre. Ogni estate erano non uno ma \emph{due} anni più grandi di me:
i gemelli avevano compiuto quattordici anni ma io ne avrei avuti ancora
dodici per altri quattro mesi frustranti. La differenza si traduceva in
vantaggio fisico. Diane sapeva benissimo che non avrei potuto batterla
sulla collina ma comunque si allontanò pedalando, e io sospirai e cercai
di pompare il mio vecchio catorcio cigolante per competere in maniera
accettabile. Non c'era gara. Diane si sollevò sul suo splendente mezzo
in alluminio inciso, e nel momento in cui raggiunse il pendio aveva
acquistato una velocità fortissima. Un trio di bambine che stava
disegnando sul marciapiede con il gesso sgambettò via per levarsi dalla
sua traiettoria. Mi lanciò un'occhiata voltandosi indietro, per metà
incoraggiante, per metà beffarda.

La strada in salita si riprese la velocità ma lei scalò i rapporti con
destrezza e rimise le gambe all'opera. Jason, ormai in cima, si era
fermato e manteneva l'equilibrio con la lunga gamba, guardando indietro
in modo interrogativo. Continuai ad arrancare, ma a metà della collina
la mia vecchia bici più che muoversi ondeggiava, e allora fui costretto
a scendere e a portarla a piedi su un fianco per il resto della salita.

Quando finalmente la raggiunsi, Diane sogghignò.

«Hai vinto» sospirai.

«Mi dispiace, Tyler. Non è stato per nulla leale.»

Scrollai le spalle, imbarazzato.

Lì la strada terminava in una via cieca, dove erano stati delineati con
pali e corde alcuni terreni residenziali ma non era stata costruita
nessuna casa. Il centro commerciale si stendeva su un lungo pendio
sabbioso verso ovest. Un sentiero di terra pressata passava attraverso
alberi intisichiti e cespugli di bacche. «Ci vediamo in fondo» esclamò,
e scomparve nuovamente.

\separatorefiori

Legammo le biciclette a una rastrelliera ed entrammo nella navata a
vetri del centro commerciale. Era un ambiente rassicurante, soprattutto
perché era cambiato pochissimo dall'ultimo ottobre. Seppure i giornali e
la televisione dessero ancora l'allarme rosso, il centro commerciale
viveva in un'aura di santa negazione. L'unica prova che qualcosa poteva
essere andato storto nel mondo esterno era l'assenza di antenne
paraboliche esposte nelle catene di negozi di elettronica e un'ondata di
pubblicazioni legate ai fatti accaduti in ottobre sugli scaffali delle
librerie. Jason soffiò su un libro in brossura con la copertina lucida
blu e oro, un libro che pretendeva di collegare l'Evento di ottobre alle
profezie bibliche: «La profezia più facile,» disse, «è quella che
predice le cose che sono già accadute.»

Diane lo guardò esasperata. «Non devi prenderlo in giro solo perché non
ci credi.»

«Tecnicamente, sto solo prendendo in giro la copertina. Non ho letto il
libro.»

«Forse dovresti.»

«E perché mai? Cosa ci sarebbe di tanto importante in questo libro per
te, da prenderne le difese?»

«Non sto difendendo niente. Ma forse Dio c'entra qualcosa con lo scorso
ottobre. Non mi pare così ridicolo.»

«In realtà,» obiettò Jason, «sì, \emph{è} ridicolo.»

Lei alzò gli occhi al cielo, fece un sospiro sonoro e si allontanò a
lunghe falcate, mollandoci sul posto. Jase rimise a posto il libro
nell'espositore e ci avviammo dietro di lei.

Gli dissi che secondo me la gente voleva solo capire cosa fosse
accaduto, ecco perché c'erano libri simili.

«O forse la gente vuole solo fingere di capire. Si chiama ``negazione''.
Vuoi sapere una cosa, Tyler?»

«Certo» risposi.

«Non lo dici a nessuno?» Abbassò la voce in modo che neanche Diane, a
pochi metri di distanza, potesse sentirlo. «Non è ancora pubblica.»

Una delle cose straordinarie di Jason era che spesso veniva a conoscenza
di cose importantissime uno o due giorni in anticipo rispetto ai
notiziari serali. In un certo senso la Rice Academy era solo la scuola
di tutti i giorni; la sua vera istruzione si svolgeva sotto la tutela di
suo padre e sin dall'inizio E.D. aveva voluto che capisse come il
commercio, la scienza e la tecnologia si intrecciano con il potere
politico. Lo stesso E.D. aveva sfruttato queste possibilità. La perdita
dei satelliti di telecomunicazione aveva aperto un nuovo grosso mercato
civile e militare per i palloni stazionari ad alta quota (gli aerostati)
che la sua azienda produceva. Si stava diffondendo una tecnologia di
nicchia e E.D. cavalcava la cresta dell'onda. E talvolta condivideva con
il figlio quattordicenne segreti che non avrebbe osato nemmeno
sussurrare a un concorrente.

E.D., ovviamente, non sapeva che Jase di tanto in tanto condivideva quei
segreti con me. Ma io stavo attento a tenerli ben stretti. (E comunque,
a chi avrei potuto raccontarli? Non avevo altri amici veri. Vivevamo in
un quartiere di gente ricca, in cui le distinzioni di classe venivano
misurate con precisione millimetrica: i figli seri e studiosi delle
madri single lavoratrici non comparivano in nessuna lista d'invitati di
serie A.)

Abbassò ancora un po' la voce. «Conosci i tre cosmonauti russi? Quelli
che erano in orbita lo scorso ottobre?»

Dispersi e presumibilmente morti la notte dell'Evento. Annuii. «Uno di
loro è vivo» disse. «È vivo e a Mosca. I russi non dicono molto, ma
corre voce che sia completamente matto.»

Rimasi a guardarlo con gli occhi sgranati ma lui non aggiunse altro.

\separatorefiori

Ci volle una decina d'anni perché fosse resa pubblica la verità ma
quando finalmente venne pubblicata (come nota in calce a una storia
europea dei primi anni dello Spin) mi tornò in mente la giornata al
centro commerciale. Ecco ciò che avvenne: tre cosmonauti russi erano in
orbita la notte dell'Evento di ottobre, di ritorno da una missione di
gestione interna alla moribonda Stazione Spaziale Internazionale. Poco
dopo la mezzanotte secondo il fuso orario della costa orientale degli
Stati Uniti, il comandante della missione, un certo colonnello Leonid
Glavin, notò la perdita di segnale dalla torre di controllo e tentò
varie volte ma senza successo di ristabilire il contatto.

La situazione, che già doveva essere allarmante per i cosmonauti,
precipitò velocemente. Quando dal Soyuz videro il Sole sorgere, ebbero
l'impressione che il pianeta attorno a cui stavano girando fosse stato
sostituito da un globo nero privo di luce.

Il colonnello Glavin alla fine lo avrebbe descritto proprio in questo
modo: un'oscurità, un'assenza visibile solo quando occludeva il Sole,
un'eclissi permanente. Il rapido ciclo orbitale dell'alba e del tramonto
era l'unica testimonianza visiva a convincerli che la Terra esisteva
ancora. La luce solare appariva all'improvviso da dietro il disco
sagomato, non rifletteva nell'oscurità sottostante e scompariva
altrettanto repentinamente quando la capsula scivolava nella notte.

I cosmonauti non potevano aver capito ciò che era successo, dovevano
essere stati terrorizzati all'inverosimile.

Dopo una settimana trascorsa in orbita nella vuota oscurità sottostante,
piuttosto che rimanere nello spazio o provare ad attraccare la vuota
ISS, i cosmonauti votarono per tentare un rientro non assistito - per
morire sulla Terra o quel che era diventata piuttosto che morire di fame
da soli. Senza la guida della torre di controllo o punti di riferimento
visivi, furono costretti ad affidarsi ai calcoli estrapolati dalla loro
ultima posizione nota. Di conseguenza, la capsula Soyuz rientrò
nell'atmosfera con un angolo pericolosamente ripido, assorbì forze
gravitazionali brutali e perse un paracadute cruciale durante la
discesa.

La capsula si schiantò con forza su una collina boscosa nella valle
della Ruhr. Vassily Golubev morì sul colpo; Valentina Kirchoff subì un
trauma cranico e morì in poche ore. Il colonnello Glavin, stordito, con
solo un polso rotto e piccole abrasioni, riuscì a uscire dal veicolo e
infine fu trovato da una squadra tedesca di ricerca e soccorso e
riaffidato alle autorità russe.

Dopo ripetuti interrogatori, i russi giunsero alla conclusione che
Glavin aveva perso la ragione a causa di quella tremenda esperienza. Il
colonnello continuava a insistere sul fatto che lui e il suo equipaggio
avevano trascorso tre settimane in orbita, ma ovviamente era pura
follia\ldots{}

Perché la capsula Soyuz, come ogni altro manufatto orbitante recuperato,
era caduta sulla Terra la stessa notte dell'Evento di ottobre.

\separatorefiori

Pranzammo nell'area ristorazione del centro commerciale, dove Diane
individuò tre ragazze della Rice che conosceva. Erano più grandi e per
me troppo sofisticate, con i capelli colorati di azzurro o rosa,
indossavano costosi pantaloni scampanati a vita bassa e catenine con
piccole croci d'oro attorno ai colli pallidi. Diane appallottolò
l'incarto dei MexiTaco e disertò il nostro tavolo per andare verso il
loro, dove tutte e quattro insieme chinarono la testa e si misero a
ridere. All'improvviso il burrito e le patatine mi divennero poco
appetitosi.

Jason si accorse della mia espressione. «Sai,» disse con delicatezza, «è
inevitabile.»

«Cosa?»

«Non vive più nel nostro mondo. Tu, io, Diane, la Grande Casa e la
Piccola Casa, il sabato al centro commerciale, la domenica al cinema.
Funzionava quando eravamo bambini. Ma non siamo più bambini.»

Non lo eravamo? No, certo che non lo eravamo, ma avevo riflettuto
seriamente su ciò che significava o avrebbe potuto significare?

«Ha il ciclo da un anno ormai» aggiunse Jason.

Sbiancai. Era più di quanto avrei dovuto sapere. Eppure, ero geloso che
lui lo sapesse e io no. Non mi aveva parlato del ciclo e neppure delle
sue amiche alla Rice. Tutte le confidenze che mi aveva fatto al
telefono, me ne resi conto d'un tratto, erano state confidenze da
bambini, storie su Jason e i suoi genitori e su quello che detestava per
cena. Ma ecco la prova, le cose che mi aveva nascosto erano tante quante
quelle che mi aveva raccontato; ecco una Diane che non avevo mai
conosciuto, che si presentava allegramente a un tavolo dall'altra parte
della navata.

«Dovremmo andare a casa» dissi a Jason.

Mi guardò commiserandomi. «Se vuoi.» Si alzò.

«Non dici a Diane che ce ne stiamo andando?»

«Penso sia occupata, Tyler. Credo abbia trovato qualcosa da fare.»

«Ma deve tornare con noi.»

«No, non deve.»

Mi sentii offeso. Non poteva scaricarci in quel modo. Non era degno di
lei. Mi alzai per andare verso il tavolo di Diane. Lei e le tre amiche
mi rivolsero tutta la loro attenzione. Concentrai lo sguardo su di lei,
ignorando le altre. «Stiamo andando a casa» la informai.

Le tre ragazze della Rice scoppiarono a ridere. Diane sorrise soltanto,
imbarazzata, e rispose: «Okay, Ty. Fantastico. Ci vediamo dopo.»

«Ma\ldots»

Ma cosa? Non mi stava neanche più guardando.

Mentre mi allontanavo, sentii una delle amiche chiederle se ero «un
altro fratello.»

«No,» rispose, «solo un ragazzino che conosco.»

\separatorefiori

Jason, che era diventato fastidiosamente sensibile, propose di
scambiarci le biciclette per tornare a casa. In quel momento non mi
interessava per niente la sua bici, ma pensai che uno scambio potesse
essere un buon modo per mascherare i miei sentimenti.

E così ce ne tornammo in cima alla Bantam Hill Road, nel punto in cui il
selciato scendeva come un nastro nero verso le strade ombreggiate dagli
alberi. Sentivo il pranzo come un blocco di cemento incastrato sotto le
costole. Esitai alla fine della via cieca, squadrando il pendio ripido
della strada.

«Scivola giù» mi incitò Jason. «Vai. Senti che sensazione si prova.»

La velocità mi avrebbe distratto? C'era qualcosa che avrebbe potuto
distrarmi? Mi odiavo per essermi permesso di credere di essere al centro
del mondo di Diane, quando in realtà ero solo un ragazzino che
conosceva.

Però era davvero una bici fantastica quella che Jason mi aveva prestato.
Mi sollevai sui pedali, sfidando la gravità a fare del suo peggio. Le
gomme raschiavano il selciato polveroso ma le catene e i cambi erano
morbidi, silenziosi, tranne che per il delicato ronzio dei cuscinetti.
Il vento scivolava dietro di me mentre acquistavo velocità. Sfrecciai
accanto a case graziosamente verniciate e ad automobili costose
parcheggiate nei vialetti, il cuore a pezzi ma libero. Quasi in fondo
alla strada cominciai a stringere i freni, perdendo slancio senza
rallentare davvero. Non volevo fermarmi. Non volevo fermarmi mai. Era
un'ottima corsa.

Ma il selciato si fece pianeggiante e finalmente frenai, mi piegai su un
fianco e riuscii a fermarmi usando la scarpa sinistra sull'asfalto. Mi
voltai.

Jason era ancora in cima alla Bantam Hill Road sopra la mia bici
sgraziata, così lontano da sembrare un cavaliere solitario in un vecchio
film western. Gli feci un cenno. Ora toccava a lui.

Sicuramente Jason aveva percorso quella collina, in salita e in discesa,
un migliaio di volte. Ma non l'aveva mai fatto su un rottame arrugginito
di seconda mano.

Lui ci stava meglio di me. Aveva le gambe più lunghe e il telaio non lo
faceva sembrare basso, ma prima di allora non ci eravamo mai scambiati
le bici; in quel momento mi vennero in mente tutti i difetti e le
particolarità della mia e quanto la conoscessi nei minimi dettagli,
pensai a come avevo imparato a non sterzare a destra troppo in fretta
perché il telaio era un po' fuori centro, a come si doveva tenere sotto
controllo l'oscillazione, a quanto fosse ridicolo il cambio. Jason non
ne sapeva niente. La collina poteva essere insidiosa. Volevo dirgli di
prendere la discesa piano ma anche se avessi urlato non mi avrebbe
sentito; mi ero allontanato troppo. Sollevò i piedi come un grosso
bambino goffo. La bicicletta era pesante. Ci volle qualche secondo per
acquistare velocità ma sapevo quanto sarebbe stato difficile fermarsi.
Era tutta massa, senza alcuna grazia. Con le mani stringevo freni
immaginari.

Credo che Jason si sia reso conto di essere nei guai solo a tre quarti
della discesa. Fu allora che la catena incrostata di ruggine si spezzò e
gli sferzò la caviglia. Ormai era abbastanza vicino da riuscire a
vederlo fare una smorfia e gridare. La bici traballò ma,
miracolosamente, riuscì a tenerla in piedi.

Un pezzo della catena si aggrovigliò nella ruota posteriore, dove si
avvolse ai montanti, producendo un suono da martello pneumatico rotto.
Due case più su, una donna che stava strappando le erbacce dal giardino
si tappò le orecchie e si voltò a guardare.

La cosa incredibile fu che Jason riuscì a mantenere il controllo del
veicolo per parecchio tempo. Jase non era un atleta ma era a suo agio in
quel suo corpo alto e smilzo. Mise fuori i piedi per tenere l'equilibrio
- i pedali erano inutili - e mantenne dritta la ruota anteriore, mentre
quella posteriore si bloccò slittando. Tenne duro.

Quello che mi stupiva era il modo in cui il suo corpo si rilassava
anziché irrigidirsi, come se fosse impegnato nel risolvere un problema
difficile ma molto interessante, come se credesse con assoluta certezza
di poter contare sulla combinazione di mente, corpo e bici per mettersi
in salvo.

Ma il mezzo lo tradì. Quel frammento di catena untuosa che sbatacchiava
pericolosamente si incastrò fra il cerchio e il telaio. La ruota, già
indebolita, si disallineò all'inverosimile e poi si piegò, sparpagliando
squarci di gomma e cuscinetti a sfera. Jason si liberò dal telaio e
rotolò in aria come un manichino gettato dall'alto di una finestra.
Batté sul selciato prima con i piedi, poi con le ginocchia, i gomiti e
la testa. Riuscì a fermarsi mentre la bici spaccata proseguì e,
superandolo, atterrò nella cunetta sul ciglio della strada, con la ruota
anteriore che continuava a girare e sferragliare. Lasciai cadere la sua
bicicletta e corsi da lui.

Si rovesciò sul dorso e alzò lo sguardo, per un attimo frastornato.
Aveva i pantaloni e la camicia strappati; la fronte e la punta del naso
erano sbucciate in profondità e sanguinavano copiosamente, la caviglia
era lacerata. Gli occhi gli si inumidirono dal dolore. «Tyler,» disse,
«oh, ahi, ahi\ldots{} mi dispiace per la tua bici, amico.»

Non che voglia dare troppo peso a quell'incidente ma ci pensai di tanto
in tanto negli anni a seguire: Jason, il suo corpo e il mezzo su cui
procedeva bloccati in un'accelerazione pericolosa e la sua composta
convinzione di poterne venire fuori bene, con le sue forze, se solo ci
avesse provato sul serio, se solo non avesse perso il controllo.

\separatorefiori

Lasciammo la mia bici ormai distrutta nella cunetta e io portai il suo
veicolo di lusso a piedi verso casa. Jason si trascinava al mio fianco,
era dolorante ma cercava di non darlo a vedere, teneva la mano destra
sulla fronte sanguinante come se avesse un brutto mal di testa, e
immaginavo lo avesse davvero.

Tornati alla Grande Casa, i suoi genitori scesero le scale del portico
per venirci incontro sul vialetto. E.D. Lawton, che doveva averci notato
dal suo studio, sembrava arrabbiato e allarmato, la bocca increspata in
una smorfia e le sopracciglia addossate agli occhi penetranti. Dietro di
lui, la mamma di Jason sembrava distante, meno interessata e forse anche
un po' ubriaca, a giudicare dal modo in cui ondeggiò quando uscì dalla
porta.

E.D. controllò le condizioni di Jase - che all'improvviso sembrava molto
più piccolo e meno sicuro di sé - e gli disse di correre in casa a
ripulirsi.

Poi si voltò verso di me.

«Tyler» disse.

«Signore?»

«Suppongo che non sia tua la responsabilità. Spero sia così.»

Aveva notato che la mia bici non c'era e che quella di Jason era
intatta? Mi stava accusando di qualcosa? Non sapevo cosa dire. Guardai
verso il prato.

E.D. sospirò. «Ti spiego una cosa. Tu sei amico di Jason. E questo è un
bene. Jason ne ha bisogno. Ma devi capire, come lo capisce tua madre,
che la tua presenza qui presuppone certe responsabilità. Se vuoi passare
del tempo con Jason, mi aspetto che tu ti prenda cura di lui. Mi aspetto
che usi il buon senso. Forse per te è un ragazzo normale. Ma non lo è.
Jason è dotato e ha un futuro davanti a sé. Non possiamo permettere che
ci siano interferenze.»

«Giusto» intervenne Carol Lawton, e in quel momento seppi per certo che
la mamma di Jason beveva. Piegò la testa e quasi inciampò sullo strato
di ghiaia che separava il vialetto dalla siepe. «Giusto, è un cazzo di
genio, è un genio, cazzo. Sarà il genio più giovane del MIT. Non
romperlo, Tyler, è fragile.»

E.D. non distolse lo sguardo. «Va' dentro, Carol» disse con voce piatta.
«Ci siamo capiti, Tyler?»

«Sissignore» risposi mentendo.

Non lo capivo proprio, E.D. Ma sapevo che una parte di quello che aveva
detto era vero. Sì, Jason era speciale. E sì, era mio dovere prendermi
cura di lui.